\documentclass[preprint,12pt]{elsarticle}

\usepackage{amsmath,amssymb,graphicx}
\journal{Internet of Things}

\begin{document}
\begin{frontmatter}

\title{Topology-aware joint RSSI forecasting for network-wide LoRaWAN monitoring}

\author[inst1]{Author One}
\author[inst1]{Author Two\corref{cor1}}
\cortext[cor1]{Corresponding author.}
\address[inst1]{Institution, city, country}

\begin{abstract}
Forecasting received signal strength can support proactive monitoring of low-power
wide-area networks, yet conventional studies commonly train one predictor for each
radio link. We formulate LoRaWAN RSSI forecasting as a network-wide multivariate task
in which one model produces forecasts for all observed nodes. The proposed
spatio-temporal graph neural network combines a geographical prior with an adaptive
adjacency learned from data, thereby representing both physical proximity and latent
cross-link dependencies. A leakage-safe protocol compares dedicated and joint
predictors on two real deployments using chronological splits, training-only scaling,
strictly consecutive hourly windows, and per-node error as the primary outcome.
Statistical baselines follow Box--Jenkins identification and include dedicated ARIMAX
and joint VARX; neural controls include dedicated and joint
sequence-to-sequence models. Model orders are selected using training information
criteria, whereas neural hyperparameters and checkpoints use validation data only.
Confirmatory results over multiple gateways and random seeds will quantify whether
joint graph forecasting improves link-level accuracy while reducing the number of
independently trained, versioned and served predictors. [TODO: insert frozen numerical
results and uncertainty intervals.]
\end{abstract}

\begin{keyword}
LoRaWAN \sep RSSI forecasting \sep graph neural network \sep multivariate time series
\sep Internet of Things \sep network monitoring
\end{keyword}
\end{frontmatter}

\section{Introduction}
LoRaWAN enables long-range connectivity for constrained Internet of Things (IoT)
devices, but its received signal strength varies with propagation conditions, weather,
and local obstruction. Prior work has shown that RSSI can be forecast with statistical
and machine-learning models and that parsimonious ARIMA models remain strong
competitors~\cite{guerra2024forecasting}. The common per-link formulation, however,
ignores dependencies among nodes received by the same gateway and requires a separate
predictor to be trained and managed for every link.

This study asks a network-level question: can a single model forecast all links at
once, retaining or improving per-node accuracy while consolidating model management?
This framing matches IoT ecosystems as multicomponent systems with dynamic
dependencies and uses graph learning to encode the monitored deployment rather than
treating its time series as unrelated channels. Graph-based models have already proved
useful for representing heterogeneous IoT structure and joint spatial--temporal
dynamics~\cite{ali2024powerful,abdelbaky2024gcn}.

The contributions are:
\begin{itemize}
  \item a network-wide formulation that predicts every LoRaWAN RSSI series with one
        jointly trained model;
  \item a physical-adaptive graph that combines distance-based structure and learned
        cross-node affinities;
  \item a paired comparison of dedicated and joint statistical and neural predictors,
        with per-node accuracy as the primary outcome and model-instance count as the
        primary consolidation measure; and
  \item a reproducible, leakage-safe evaluation on distinct real deployments, with
        Box--Jenkins identification and explicit integration/cointegration checks.
\end{itemize}

\section{Related work}
\subsection{RSSI forecasting in LoRaWAN}
Guerra et al.~\cite{guerra2024forecasting} studied weather-assisted LoRaWAN RSSI
forecasting using ARIMA, machine learning, LSTM and hybrid models. Our work changes the
unit of modelling from one link to the gateway-level system and evaluates whether
cross-link information can improve forecasts while consolidating predictors. Reviews
of LoRaWAN optimization also emphasize that network decisions span interacting layers
and operating conditions~\cite{loraReview2024}.

\subsection{Joint spatio-temporal learning in IoT}
Graph neural networks provide an inductive bias for variables embedded in a known or
partly known topology. Physical structure alone can be incomplete, motivating learned
adjacency mechanisms such as Graph WaveNet and multivariate graph learning
networks~\cite{wu2019graph,wulog2020connecting}. Recent work in \emph{Internet of
Things} combines graph convolution and recurrent modelling for dynamic IoT links and
demonstrates the journal's interest in graph representations of IoT
systems~\cite{abdelbaky2024gcn,ali2024powerful}.

\section{Problem formulation}
Let $\mathbf{x}_t\in\mathbb{R}^{N}$ contain the hourly mean RSSI of $N$ nodes observed
by one gateway, and let $\mathbf{w}_t$ denote historical weather covariates. Given a
history of length $L$, the task is
\begin{equation}
 f_\theta:\{\mathbf{x}_{t-L+1:t},\mathbf{w}_{t-L+1:t}\}
 \mapsto \widehat{\mathbf{x}}_{t+1:t+H}.
\end{equation}
A dedicated strategy estimates $N$ independently parameterized functions; a joint
strategy estimates one function whose output contains all nodes. Performance is
reported for each node before any macro average.

\section{Methodology}
\subsection{Data control and forecasting protocol}
Each gateway is treated as a separate deployment experiment. Observations are ordered
chronologically and divided into 70\% training, 15\% validation and 15\% test
partitions. All scalers are fitted on training data only. A sample is eligible only
when the complete input and target interval is observed at exact hourly spacing for
all included nodes; missing values and temporal gaps are never imputed across. Future
weather is unavailable. For horizon $H$, statistical exogenous regressors are delayed
by $H$ hours, so every model uses only information available at the forecast origin.

The test partition is evaluated once after all decisions are frozen. Neural search and
checkpoint selection use validation macro-MAE. Statistical identification is confined
to training data and validation error is descriptive rather than selective.

\subsection{Box--Jenkins identification and multivariate diagnostics}
For each RSSI series, ADF and KPSS tests are applied in levels and first differences.
They restrict the admissible differencing order. When ADF rejects a unit root but KPSS
also rejects level stationarity, we retain $d=0$ to avoid overdifferencing and record
the disagreement as a structural-instability warning; the multivariate rank diagnostic
provides an additional system-level check. ACF and PACF are then computed
on the corresponding transformed series. A pre-specified parsimonious shortlist
contains the suggested $(p,q)$ order, its immediate neighbours, and the simple AR(1),
MA(1), and ARMA(1,1) forms; this avoids an indiscriminate grid search while retaining
diagnostic traceability. Each shortlisted ARIMAX $(p,d,q)$ is estimated by exact Gaussian state-space
likelihood over independent contiguous training segments and selected by corrected
Akaike information criterion,
\begin{equation}
 \mathrm{AICc}= -2\ell+2k+\frac{2k(k+1)}{n-k-1}.
\end{equation}
Ljung--Box statistics diagnose residual autocorrelation. The procedure produces a
separate order for every dedicated link, consistent with the Box--Jenkins principle
that model structure follows the observed series~\cite{box2015time,hyndman2008automatic}.

For joint VAR models, integration is assessed for every endogenous and exogenous
series and Johansen's trace test is applied to the RSSI vector on training data. A
full-rank conclusion, checked for one to three lagged differences, supports a
stationary level VAR. If the variables are $I(1)$ and the rank is reduced, level
VARX is not fitted: the pre-specified remedy is a
VECM/VECMX model, because differencing alone would discard the error-correction
relation~\cite{johansen1991estimation,lutkepohl2005new}. VAR lag order and regularization
are selected by training BIC. Candidate structures include dense lags $1,\ldots,p$
up to 24 hours and parsimonious seasonal sets $\{1,6,12,24\}$ and
$\{1,2,3,6,12,24\}$; unstable autoregressive dynamics are rejected. BIC is never
compared across level and differenced dependent variables.

\subsection{Neural baselines}
The dedicated Seq2Seq baseline fits one recurrent encoder--decoder per node. Its joint
counterpart encodes all RSSI and historical weather channels, then uses an
autoregressive LSTM decoder with additive attention and a last-observation residual.
Both use identical chronological partitions and stopping rules. [TODO: document the
frozen search space and selected configurations in a supplementary table.]

\subsection{Physical-adaptive spatio-temporal graph network}
The graph model represents nodes with a fixed geographical adjacency
$\mathbf{A}_{p}$ and a learned adaptive adjacency $\mathbf{A}_{a}$. The physical
component is obtained from pairwise node distances and the adaptive component from
trainable node embeddings. Their normalized combination drives graph message passing,
while causal dilated temporal convolutions encode RSSI history. Self-edges are
excluded. Ablations remove the graph, retain only the physical component, or retain
only the adaptive component. These weights are predictive associations and are not
interpreted as physical connectivity or causality.

\section{Experimental design}
\subsection{Deployments and gateways}
[TODO: add a table containing deployment, gateway, date range, node count, missingness,
and eligible train/validation/test origins. Explain gateway and node inclusion criteria
before presenting performance. The vineyard deployment contains eight nodes; the UVA
external deployment is evaluated gateway by gateway rather than pooling gateways.]

\subsection{Outcomes and statistical analysis}
Primary outcomes are MAE and RMSE in dB for every node. Macro averages, terminal-lead
and lead-wise errors are secondary. Consolidation is measured primarily by the number
of independently parameterized model instances required to cover all links, supported
by coefficient count and serialized coefficient storage. Hardware-sensitive inference
latency is not a primary scientific outcome.

Confirmatory comparisons use identical forecast origins. [TODO: pre-specify repeated
seeds for neural models, dependence-aware confidence intervals (for example a moving
block bootstrap over forecast origins), paired effect sizes, multiplicity handling,
and the rule for aggregating evidence across gateways.]

\section{Results}
\subsection{Stationarity, integration and model identification}
[TODO: report ADF/KPSS conclusions, Johansen rank by gateway, selected ARIMAX orders,
VAR orders, convergence and Ljung--Box diagnostics. Keep full numerical diagnostics in
the supplement/repository.]

\subsection{Per-node forecasting accuracy}
[TODO: lead with a model-by-node MAE heatmap/table; report uncertainty and paired
differences. Global means are supplementary and must not replace node-level results.]

\subsection{Graph ablations and operational consolidation}
[TODO: compare hybrid, physical-only, adaptive-only and no-graph variants. Report
instances, parameter count and storage without claiming unmeasured energy or latency
benefits.]

\section{Discussion}
The central interpretation is not that graph learning universally dominates all
forecasters, but whether shared cross-link structure yields a consistent accuracy and
management trade-off across gateways. Improvements of the adaptive-only model would
support latent cross-link dependence; an additional gain from the hybrid graph would
support the geographical prior. Gateway-specific failures and nodes that do not
benefit are equally important to the conclusion.

Hourly mean RSSI forecasting does not directly establish packet-level ADR, delivery
ratio, energy consumption or outage prevention. Complete-window evaluation also does
not demonstrate robustness to nodes disappearing online. These are downstream system
questions rather than outcomes measured here.

\section{Conclusion}
[TODO: write after freezing confirmatory experiments. State the network-wide finding,
its cross-deployment consistency, the consolidation trade-off, and no stronger claim.]

\section*{Data and code availability}
[TODO: add repository release, commit, licenses and permanent dataset identifiers.]

\section*{CRediT authorship contribution statement}
[TODO: list author roles.]

\section*{Declaration of competing interest}
The authors declare that they have no known competing financial interests or personal
relationships that could have appeared to influence the work reported in this paper.

\bibliographystyle{elsarticle-num}
\bibliography{references}
\end{document}
